% 01-introduction — v4.0
\subsection*{A device in search of a method}
Low-cost EEG headsets are sold as attention trainers, relaxation aids, sleep gadgets and
emotion sensors. Some of them record real neural signal: consumer devices have been
validated against research amplifiers for event-related potentials
\citep{krigolson2017muse,badcock2015emotiv}, and resting-state comparisons show that
spectral estimates and test--retest behaviour differ across systems
\citep{ratti2017comparison,cannard2021muse}. The problem is what happens after the
recording. A paper may report that alpha power increased during relaxation without
saying which electrode was used, how the signal was filtered, whether blinks were
removed, or how the result compares with a conventional montage. A reader cannot tell
whether the finding reflects a cognitive process or an artefact of the recording setup.

This matters for two reasons. First, the same claim can be produced by very different
pipelines, so the result is difficult to reproduce. Second, a single frontal electrode
picks up cortical activity together with eye movements and muscle noise
\citep{whitham2007emg}; presenting that mixture as a specific cognitive measure invites
overclaiming. Device-comparison studies are useful, but they answer a different
question: whether one headset is close enough to a reference amplifier under one
protocol. They do not tell a reader which analyses still work when the acquisition is
reduced.

\subsection*{What this paper does}
We take high-quality open recordings and degrade them one step at a time towards what a
cheap headset provides. The steps are explicit: a single electrode, 128\,Hz sampling,
1\,\textmu V quantisation, 1--40\,Hz band-limiting, pink noise at four levels, and
synthetic blinks at three rates. Each step is reported separately, so a reader can see
which manipulation matters. The benchmark has four analysis families: alpha
reactivity, affective decoding, motor decoding (exploratory) and reliability. The
reference is always the same recording before degradation, and, where it matters, the
same electrode rather than a whole-head average.

Two open datasets make this feasible. The EEG Motor Movement/Imagery Dataset provides
109 subjects, 64 channels and eyes-open and eyes-closed baselines
\citep{schalk2009eegbci}. An open affective dataset contains 20 participants recorded
simultaneously with a low-cost Emotiv EPOC+ and a high-end BrainVision system, with
valence--arousal quadrant labels \citep{sanchezreolid2022emotion,sanchezreolid2025dataset}.
The first dataset is used for the degradation grid; the second lets us compare a real
low-cost device with a high-end montage and with the same high-end data reduced to one
electrode.

\subsection*{What we find}
The results are mixed, and the mix is the point. Alpha reactivity is robust to the
acquisition chain: once the reference is taken at the same electrode, 99.99\% of the
effect remains. The choice of electrode, however, costs about 42\% of eyes-closed alpha
power, and noise reduces the effect monotonically. In the affective data, the full
montage decodes valence, arousal and quadrant above chance, but a single electrode does
not: valence is no longer detectable, arousal weakens, and four-class accuracy drops
close to the level expected from a weak classifier. A band control suggests that part of
the montage advantage comes from beta-band features, which may include muscle activity.
A 14-channel Emotiv performs about as well as the high-end montage in this sample, but
the equivalence test does not pass, so the comparison remains descriptive.

The contributions are: a reproducible degradation pipeline; a decision matrix that maps
each result to a claim-level verdict; a same-site retention estimate for alpha; a
decomposition of the affective result into representation, site and band effects; and a
reporting checklist for single-channel studies, based on COBIDAS MEEG
\citep{pernet2020cobidas}, EEG-BIDS \citep{pernet2019bids} and the CRED-nf checklist
\citep{ros2020crednf}. The code, the derived results and the JSON files behind every
number are archived with the paper.

The paper is not a clinical or wellness intervention study, and it does not validate a
commercial product. It is a methods study: it measures what different analyses can
support, and it shows where the limits lie.
